\uuid{aWHp}
\exo7id{6542}
\titre{exo7 6542}
\auteur{drutu}
\organisation{exo7}
\datecreate{2011-10-16}
\isIndication{false}
\isCorrection{false}
\chapitre{Polynôme}
\sousChapitre{Polynôme}
\cpge{mpsi-mp2i.structures-algebriques.c, mpsi-mp2i.polynomes-fractions-rationnelles.c, pcsi.polynomes.c, ptsi.polynomes.c, mp-mpi.structures-algebriques.b, mp-mpi.structures-algebriques.d}
\module{Algèbre et théorie des nombres}
\niveau{L3}
\difficulte{}

\contenu{
\texte{
Soit $A$ un anneau commutatif. Montrer l'équivalence des deux affirmations suivantes :
}
\begin{enumerate}
    \item \question{[(i)] $A $ est un corps fini ;}
    \item \question{[(ii)] tout polyn\^ome $P\in A[X]$ de degré $n\geq 1$ a au plus $n$ racines dans $A$ et toute fonction $f:A \to A$ est polynomiale.}
\end{enumerate}
}
